\section{Discussion}
\label{sec:discussion}

\subsection{From Single-Day to Multi-Day: Bridging Two Temporal Scales}

The progression from single-day detection (71.5\%) to multi-day regime identification (81.2\% in 2024) reveals complementary validation dimensions. Single-day testing establishes that the LLM understands dealer mechanics within individual trading sessions; regime detection validates that this understanding extends to persistent multi-day structural states. The raw chain finding (92.3\% detection from first-principles data) provides additional support: if the model merely matched statistical signatures, raw chain performance would degrade rather than improve when pre-calculated summaries are removed.

The detection-alpha orthogonality finding from single-day analysis (Sharpe collapse from 1.8 to 0.1 while detection remains stable) motivated the regime extension: patterns that persist despite economic insignificance must reflect structural mechanics rather than exploitable inefficiencies. This distinction---mechanism identification versus alpha generation---is fundamental to deploying LLMs for market surveillance and risk management rather than trading.

\subsection{Validating Structural Reasoning}

Five-phase validation provides evidence that the LLM performs structural reasoning rather than pattern matching. The 69.1 percentage point discrimination between 2024 (81.2\%) and 2020 (12.1\%) establishes selectivity: if the model relied on memorized associations or universal statistical patterns, detection rates would converge across periods rather than exhibiting 6.7$\times$ separation.

Three findings support this interpretation. First, the LLM maintained 98\% mechanical accuracy across both detection conditions---correctly computing persistence, magnitude, and flip counts regardless of whether windows qualified as regimes. Second, negative controls achieved 0\% false positives on transitional and low-magnitude synthetic data, confirming the model enforces structural criteria rather than over-detecting patterns. Third, confidence scores provide continuous quality discrimination beyond binary classification ($r = +0.719$ with persistence, $r = -0.705$ with sign flips), suggesting the model tracks regime robustness, not just threshold crossing.

This selectivity contrasts with earlier 5-day trajectory testing \citep{regan2025obfuscation}, which achieved 98--100\% detection across all market conditions by identifying universal daily hedging flows. The deliberate shift to 30-day windows with strict criteria transformed the task from universal pattern matching to selective regime identification---a critical distinction for validating genuine structural reasoning.

\subsection{Market Structure Evolution and 0DTE Hypothesis}
\label{sec:discussion:0dte}

The multi-year temporal analysis (Phase 5) reveals a gradual, non-monotonic evolution in regime detection that \textit{coincides with} the adoption curve of zero-days-to-expiration (0DTE) options rather than a sharp structural break. Detection progressed from 12.2\% (2020) through borderline years (3.7\% in 2021, 32.4\% in 2022) to sustained 100\% detection in 2024--2025, with average GEX magnitude growing from \$3.0B to \$20.3B. We are careful to frame the 0DTE correspondence as temporal coincidence supported by a plausible mechanical channel (pinned daily dealer hedging demand), not as a demonstrated causal relationship; the observational design here cannot rule out alternative drivers, and we list those explicitly in \S\ref{sec:limitations}.

The non-monotonic pattern (32.4\% in 2022 dipping to 20.2\% in 2023 before reaching 100\% in 2024) is suggestive rather than definitive: 2023's elevated volatility (FOMC uncertainty, banking stress) appears to have disrupted regime formation despite growing 0DTE hedging pressure, which is consistent with the view that regime persistence requires both sustained dealer pressure \textit{and} a volatility environment permitting consolidation. We read this dynamic as consistent with, rather than proof of, a 0DTE-mediated structural interpretation; stronger causal identification would require a natural experiment---for example a temporary 0DTE suspension, a 0DTE launch on a comparable non-SPY underlier that could serve as a counterfactual, or an instrumental-variable design that separates the 0DTE channel from contemporaneous shifts.

Concurrent factors that cannot be excluded in the observational data include: (i) the 2021--2023 interest-rate cycle ($\approx$0.25\% $\rightarrow$ 5.5\%), which reshaped the carry landscape for every option-writing strategy, (ii) the growth of systematic short-volatility flow into single-name and index options, (iii) increasing passive and index-linked assets under management, and (iv) changes in market-maker concentration following the 2020--2022 period of retail-driven volatility. Each of these could plausibly contribute to the detection pattern alongside 0DTE adoption. The non-monotonic detection trajectory is \textit{less easily reconciled with} a simple gradual secular trend---detection remained low through 2023 despite continuous interest-rate increases and continuing technology diffusion, then jumped discontinuously in 2024 coincident with 0DTE market saturation ($\approx$46\% SPY volume share)---but we emphasise that ``less easily reconciled'' is not ``ruled out.'' Disentangling these channels is beyond the scope of this work, which focuses on LLM-validation methodology rather than on causal microstructure inference.

\subsection{Temporal Obfuscation as Generalizable Methodology}

Our obfuscation testing framework generalizes beyond financial applications. The core principle---stripping temporal and contextual identifiers to force structural reasoning---addresses a universal challenge in deploying LLMs for domain-specific analysis: distinguishing genuine reasoning from training data contamination. The methodology requires three components: (1) identifiable features that can be systematically removed, (2) structural patterns that persist after obfuscation, and (3) negative controls validating discrimination.

Potential applications include medical time series analysis (removing patient identifiers and dates while preserving physiological patterns), climate data analysis (obfuscating location and temporal markers), and industrial anomaly detection (stripping equipment identifiers). In each case, temporal obfuscation could validate whether LLMs reason from structural properties or rely on memorized associations from training data.

\subsection{Dispersed Knowledge and Information Aggregation}

The raw chain superiority finding (92.3\% vs.\ 61.5\%) admits a theoretical reading grounded in \citet{hayek1945}. Hayek argued that the economic problem is one of utilising knowledge ``not given to anyone in its totality''---relevant information is distributed across many participants, each holding local knowledge that resists centralised aggregation. Traditional parametric metrics like net GEX attempt exactly that centralised aggregation: they compress the full strike-level options surface into a single scalar. Our results suggest the compression is lossy in a predictable way.

The raw options chain preserves what Hayek termed ``knowledge of the particular circumstances of time and place''---strike-by-strike asymmetries, implied-volatility skew patterns, and open-interest distributions that encode actual dealer positioning. When the LLM reads the raw chain, it acts as an aggregator that synthesises these local signals into a single structural assessment. The 30.8 percentage point improvement over GEX-assisted detection is direct evidence that the ``local knowledge'' embedded in individual strikes carries structural signal that scalar aggregation destroys.

This interpretation extends beyond methodological curiosity. The alpha disappearance puzzle---stable detection (68--74\%) despite collapsing profitability (Sharpe 1.8 $\rightarrow$ 0.1)---parallels a core insight from Austrian capital theory: that structural knowledge of market coordination patterns differs fundamentally from actionable entrepreneurial knowledge \citep{kirzner1973competition}. Detecting that dealers are constrained to pro-cyclical hedging is structural knowledge; profiting from that detection requires entrepreneurial judgment about timing, magnitude, and competing market forces that our framework deliberately excludes. The orthogonality between detection and alpha is thus not a limitation but a feature: it confirms the framework identifies genuine structural mechanics rather than ephemeral trading opportunities.

\subsection{Practical Implications}
\label{sec:discussion:practical}

The results carry concrete implications along three axes that matter
most to financial practitioners: risk management, market efficiency,
and the design of quantitative research pipelines. We address each in
turn.

\subsubsection{Risk management}

Our dealer-gamma regime detection is best understood as a
\textit{risk-regime indicator} rather than a trading signal. Because
the framework achieves stable detection (68--74\% quarterly) even as
Sharpe ratios collapse from 1.8 to 0.1, its output is a reading of
whether the market is currently operating in a mechanically constrained
state---not a forecast of directional return. Three specific risk-management
applications follow.

First, for \textbf{intraday volatility budgeting}, a persistent negative
gamma regime implies amplified dealer chase-hedging and elevated realised
volatility clustering on high-volume days; sell-side risk desks can use
the 30-day regime classification as a leading indicator for
volatility-of-volatility exposure and size gamma-exposure limits
accordingly. Second, for \textbf{option-book hedging under OpEx
concentration}, the known pinning dynamic around large-OI strikes
\citep{ni2005stock} is substantially more forceful under positive-gamma
regimes that we now classify explicitly; book hedgers may increase
hedging frequency around OpEx when the framework detects a persistent
positive regime. Third, for \textbf{risk-scenario design}, 30-day
regime history provides a natural conditioning variable for stress-test
calibration: a 2020-style fragmented period and a 2024-style persistent
negative period imply materially different joint distributions for
realised volatility and spot-vol covariance.

\subsubsection{Market efficiency}

The detection-alpha orthogonality result---stable structural detection
that does not translate into exploitable profit---contributes to the
ongoing debate about efficiency in optionized equity markets. Our
evidence is consistent with a weakly efficient market in which
\textit{structural constraints} are reliably identifiable but already
priced: arbitrageurs compete away the first-order profit from knowing a
regime exists, yet the underlying mechanics that generate volatility
clustering and OpEx pinning persist because they are mandatory, not
opportunistic, actions by dealers. This reconciles two claims that are
often treated as contradictory: that dealer-gamma positioning demonstrably
influences short-horizon price dynamics \citep{anderegg2022impact}, and
that systematic strategies exploiting that influence deteriorate as
attention accumulates. The framework thus provides a positive account of
why microstructure-aware research can be genuinely informative for risk
without being genuinely informative for alpha.

\subsubsection{Practitioners: data-pipeline and model-deployment design}

Two practitioner-facing design implications follow directly from the
experimental results. First, the 30.8-percentage-point advantage of raw
strike-level data over pre-aggregated GEX (92.3\% versus 61.5\% at the
single-day scale) is evidence that common pipeline designs---which
compress option chains into scalar summaries before analysis---are
discarding signal that an LLM can reconstruct when given the raw input.
This challenges the default of parametric aggregation and suggests that
practitioners deploying LLM-based analysis should ingest granular data
wherever feasible. The design principle generalises beyond options
dealer flow to any domain where scalar metrics are conventional summary
of distributional inputs: credit risk (single probabilities from granular
exposure tapes), fixed-income surveillance (duration summaries from
granular curve positions), and equity factor research (single factor
loadings from granular return-attribution inputs) are the obvious
candidates.

Second, the 0DTE-driven regime shift detected in 2022--2024 implies
that static microstructure models calibrated to pre-2022 data will miss
a structural reorganisation that our framework picks up. Practitioners
running surveillance, risk, or execution models should treat the
2022--2024 period as a regime change requiring recalibration, not a
drift-along-the-same-curve.

\subsection{Limitations and Future Work}
\label{sec:limitations}

Seven limitations merit explicit discussion, and we describe the specific
future work each motivates.

\noindent\textbf{Single-asset scope}: All reported results concern SPY.
SPY is deliberately chosen as the highest-liquidity and earliest 0DTE-enabled
U.S.\ equity benchmark, but this choice leaves cross-asset generalization
empirically untested. Dealer positioning in QQQ, IWM, single-name equities,
and non-equity underliers (futures, rates, FX) may exhibit different
regime dynamics because of differences in option chain depth, 0DTE
availability, and the composition of end-users. Cross-asset replication is
the single highest-priority item for future work; a pre-registered protocol
applying the same obfuscation and regime-classification framework to at
least one additional ETF (QQQ) and one individual equity (e.g., NVDA, AAPL)
would directly test the transferability claim.

\noindent\textbf{Single-LLM dependence}: All 2{,}221 evaluations were
produced by a single reasoning model (OpenAI o4-mini). The detection
rates reported here are therefore conditional on this specific model's
prior distribution over market-structure reasoning. Model-swap
validation with alternative reasoning families
(e.g., Anthropic Claude, OpenAI o3, Google Gemini, open-source
reasoning models) using identical prompts and the same obfuscated
sequences is a direct and low-cost extension. A cross-model agreement
analysis would sharpen the distinction between the framework's
structural-reasoning claim and any o4-mini-specific artefacts.

\noindent\textbf{Lack of independent external validation}: Our per-window
ground-truth metrics (persistence, magnitude, sign flips) are computed
from the same Alpha Vantage options feed used to construct the windows.
We do not cross-validate detected regimes against an independent
data source (CBOE DataShop, OPRA consolidated feed, or a commercial
vendor such as SpotGamma or MenthorQ) or against an independent oracle
of dealer positioning. External validation---both against a second
options-data pipeline and against related microstructure observables
(realised volatility, implied-realised spread, opening auction
imbalance)---would strengthen the claim that the detected regimes
correspond to a real, cross-verified phenomenon rather than an artefact
of any single data provider.

\noindent\textbf{End-of-day measurement}: Our GEX\_OI approach captures
dealer inventory at the close but not intraday gamma dynamics; high-frequency
flow data could refine detection, particularly for 0DTE contracts that
expire within a single trading session. Intraday GEX surface reconstruction
from streaming OPRA is a natural extension.

\noindent\textbf{Causal attribution}: The 0DTE hypothesis is supported by
temporal coincidence and a theoretical mechanism but remains
circumstantial; observational data cannot exclude alternative explanations
such as post-pandemic monetary-policy shifts, passive-flow concentration,
or market-maker inventory changes. A natural experiment---for instance a
temporary 0DTE suspension or a regulatory halt during a market stress
episode---would provide stronger causal evidence. We treat the 0DTE
correspondence as consistent with our structural-regime detection rather
than as a demonstrated causal channel (see \S\ref{sec:discussion:0dte}).

\noindent\textbf{Shuffle test asymmetry}: The 61\% false positive rate on
2024 shuffled data (versus 12.1\% on 2020 shuffled data) reflects extreme
regime persistence rather than framework failure---2024 regimes exhibit
such dominant same-sign positioning that randomizing day order rarely
disrupts the aggregate signal. This asymmetry is itself informative,
confirming that 2024 regimes are defined by aggregate dominance rather
than temporal sequencing, but it means the shuffle test's diagnostic
power is lower in high-persistence regimes and should be interpreted
accordingly.

\noindent\textbf{Threshold sensitivity}: All tested parameter
configurations maintained substantial 2024-versus-2020 discrimination
(see \S\ref{sec:regime} sensitivity analysis), but the chosen thresholds
(persistence $\geq$70\%, magnitude $\geq$\$5B, flips $\leq$5) represent
empirically validated design choices rather than first-principles
derivations. Future work should explore adaptive thresholding that
responds to volatility regime, contract-maturity mix, or prevailing
options notional.
